\section{Discussion}

\subsection{Strategic Style Versus Raw Leaderboard Position}

The benchmark surfaces stable strategic profiles rather than only a leaderboard. Across the frozen 600-game cohort, Kimi claims 225 wins for a 37.5\% overall win share while maintaining only 11.3\% challenge frequency, whereas Mistral remains at 0 wins with 63.4\% challenge frequency. Kimi is the most robust finisher because it leads Experiments~1, 2, and~3 by win share. Nemotron is the strongest condition-specific finisher in the low-strategy control, where it takes 49 of 150 wins for a 32.7\% share. Qwen is the most consistent runner-up across the deception-relevant conditions. These are not interchangeable styles: the benchmark is exposing different ways to trade off bluffing, inspection, and risk.

\begin{figure}[H]
  \centering
  \includegraphics[width=0.96\linewidth]{lie_frequency_vs_win_rate.png}
  \caption{Lie frequency versus win share across all 24 model-condition rows. Winning is not a monotonic function of deception alone: Kimi combines strong win shares with relatively low challenge intensity, Nemotron is strongest mainly in the control condition, and Mistral occupies the high-suspicion, zero-win corner.}
  \label{fig:lie-frequency-vs-win-rate}
\end{figure}

\begin{figure}[H]
  \centering
  \includegraphics[width=0.96\linewidth]{challenge_frequency_vs_win_rate.png}
  \caption{Challenge frequency versus win share across all 24 model-condition rows. The strongest outcomes cluster away from the extreme high-challenge region, while Mistral anchors the costly over-enforcement corner.}
  \label{fig:challenge-frequency-vs-win-rate}
\end{figure}

\subsection{Instruction Compliance Under Competitive Incentives}

The honesty-mandate condition is the strongest alignment result in the paper, but only after separating optional dishonesty from conflict turns. The plain-language rule sharply suppresses lying when truthful play is available: mean model-level optional-lie rate falls from 13.4\% in Experiment~1 to 2.2\% in Experiment~3. But about a quarter of Experiment~3 turns are truthful-play-unavailable at the moment of action, so overall lie frequency still lands at 27.2\%. The table also becomes less vigilant. Relative to Experiment~1, mean challenge frequency falls from 45.3\% to 21.8\%, while mean lie success rises from 11.0\% (95\% CI 7.5--16.0) to 34.7\% (95\% CI 27.6--44.5). In this benchmark and pilot cohort, prompt framing therefore changes the social equilibrium of the table, not just the actor's willingness to bluff.

\subsection{Bullshit as an Inspection Game}

Bullshit can be read as an inspection game. If players challenge too often, they punish themselves; if they challenge too rarely, bluffing becomes cheap. The frozen cohort supports this interpretation. Experiment~1 produces the highest challenge rate and the lowest lie success, while Experiment~3 produces much lower challenge intensity and much higher lie success. A secondary observation sharpens the point: Experiments~0 and~3 converge to nearly the same enforcement regime, with mean challenge frequencies of 21.5\% and 21.8\% and mean lie-success rates of 37.9\% and 34.7\% respectively. Figures~\ref{fig:lie-frequency-vs-win-rate} and~\ref{fig:challenge-frequency-vs-win-rate} make the same calibration point from two directions: winning is not a simple function of more bluffing or more challenging. In this sense, BS-Bench measures not only who lies, but how agents calibrate signaling, inspection, and punishment under uncertainty.

The model-level Experiment~3 shifts suggest a narrower follow-on interpretation that is consistent with a role for challenge calibration in addition to optional-lie reduction. Table~\ref{tab:cross-experiment-summary} shows that mean model-level challenge accuracy does not uniformly improve across conditions, so we do not claim a condition-wide accuracy gain. The narrower point is model-level, and Table~\ref{tab:exp1-exp3-model-comparison} makes those six-model shifts visible in the main results. Qwen is the clearest positive case: its win share rises from 16.7\% in Experiment~1 to 34.0\% in Experiment~3 while its challenge frequency falls from 41.9\% to 6.1\%, its challenge accuracy rises from 43.6\% to 55.7\%, and its adjusted optional-lie rate falls from 15.3\% to 0.4\%. MiniMax is the counterexample that keeps the claim narrow: its adjusted optional-lie rate also falls sharply, from 7.2\% to 0.8\%, but its challenge accuracy dips from 38.6\% to 35.7\% and its win share falls from 16.7\% to 4.0\%. Mistral remains the negative anchor: it posts the highest Experiment~3 challenge frequency at 50.8\%, the highest adjusted optional-lie rate at 11.7\%, and still wins 0.0\% of games. These comparisons are descriptive rather than causal, and they suggest only that some Experiment~3 winners combine lower challenge frequency with better-targeted challenges.

Mistral is the clearest example of costly over-enforcement. It challenges more than every other model in every condition---62.5\% in Experiment~0, 74.7\% in Experiment~1, 71.3\% in Experiment~2, and 50.8\% in Experiment~3---yet it never wins a condition. This is evidence for a narrow game-theoretic claim about this benchmark: indiscriminate suspicion is a losing strategy. Challenge behavior matters, but only when it is calibrated well enough that the punisher is not paying the cost of constant failed enforcement.

The turn-level logs also show that some models attempt much stronger bluffs than the low-exposure examples highlighted in Section~4. Under a strict definition---claiming four cards of a rank while actually playing zero matching cards---the control condition contains only 10 such bluffs and none survive. The prompting conditions are materially higher: Experiment~1 contains 1{,}369 such bluffs with 27 unchallenged, Experiment~2 contains 997 with 15 unchallenged, and Experiment~3 still contains 277 with 10 unchallenged. The survival rates are therefore 0.0\% in the control, 2.0\% in Experiment~1, 1.5\% in Experiment~2, and 3.6\% in Experiment~3. The descriptive point is twofold. Prompting changes how often these egregious bluffs are attempted, and it also changes whether the table lets any of them survive. Under a broader definition of any three-or-four-card zero-match claim, the frozen cohort contains 3{,}259 such bluffs, with 78 unchallenged. So the benchmark does not only contain cautious one-card embellishments. It also contains egregious zero-match bluffs, and a small fraction of them survive.

The reasoning traces are consistent with this interpretation. In one Experiment~3 game, MiniMax explicitly notes that it has a truthful 6 available and still pads the play with a 5H to make a two-card 6-claim harder to challenge. In another Experiment~3 game, Qwen acknowledges the honesty rule but also that it has no legal truthful play of the required rank, making the resulting bluff better interpreted as a conflict turn than as an optional lie. A third trace shows Nemotron challenging an impossible two-card claim by noting that it already holds three of the four relevant cards. Short excerpts are included in Appendix~\ref{app:reasoning-traces} to give a concrete sense of how these strategic choices appear in the logs.

\subsection{Limitations}

The current paper studies one environment family, one primary provider cohort, and one roster of six hosted models. Experiment~0 should be read as a low-strategy prompt control, not a pure non-strategic baseline, because the shared rules block still exposes lying, challenging, and risk language; Exp0-to-Exp1 differences are therefore best read as changes in prompt intensity rather than as strategy appearing from zero. The frozen cohort preserves actor reasoning and actual challenge rationales, but not comparable pass rationales from players who decline to challenge, so interpretations of unchallenged bluffs rely on observed behavior rather than direct evidence of why judges passed. The paper does not include a human baseline in the main pilot, and it does not support claims about model size effects or broad real-world deception beyond the benchmark. A stronger v2 should add a neutered-rules control with no deception, challenge, or risk language in the shared prompt, plus human player baselines to calibrate absolute behavior levels. The pilot also does not establish internal motives or justify sweeping generalization from one card game to all strategic settings. These are limitations of scope, not reasons the benchmark is invalid.
